% ============================================================
% semana_03_conjuntos — ejercicios
% ============================================================
% Este archivo es incluido desde main.tex con \include{}

% ╔══════════════════════════════════════════════════════════════════════════╗
% ║  EJERCICIOS PROPUESTOS — TEORÍA DE CONJUNTOS                             ║
% ║  Curso: Aritmética | Semana: 01                                           ║
% ╚══════════════════════════════════════════════════════════════════════════╝

\newpage
\section{Ejercicios propuestos --- Teor\'{\i}a de Conjuntos}

\begin{tcolorbox}[
  enhanced, colback=GrisClaro, colframe=AzulPizarra,
  arc=2pt, boxrule=0.7pt, left=8pt, right=8pt, top=4pt, bottom=4pt
]
\small\sffamily
\textbf{Instrucciones:} Marca la alternativa correcta. Cada ejercicio tiene
una \'unica respuesta v\'alida. Tiempo estimado: \textbf{30 minutos}.
\hfill \textbf{Semana 01 --- Teor\'{\i}a de Conjuntos}
\end{tcolorbox}

\vspace{0.4cm}

% ─── NIVEL BÁSICO ────────────────────────────────────────────────────────────
\begin{tcolorbox}[
  enhanced, colback=VerdeSuave!40, colframe=VerdeOliva,
  arc=2pt, boxrule=0.5pt, fonttitle=\sffamily\bfseries\small,
  title={Nivel B\'asico}, left=4pt, right=4pt, top=3pt, bottom=3pt
]\end{tcolorbox}

\begin{multicols}{2}
\setcounter{numejercicio}{0}

\ejercicio Dado el conjunto $A = \{1,\; 3,\; \{4\},\; 8\}$,
¿cu\'al de las siguientes proposiciones es \textbf{verdadera}?
\begin{alternativas}
  \item $4 \in A$
  \item $\{4\} \notin A$
  \item $\{8\} \in A$
  \item $3 \notin A$
  \item $\{4\} \in A$
\end{alternativas}

\vspace{0.3cm}

\ejercicio Si $B = \{2,\; 2,\; 4,\; 5,\; 5,\; 5\}$, entonces $n(B)$ es:
\begin{alternativas}
  \item $6$
  \item $5$
  \item $4$
  \item $3$
  \item $2$
\end{alternativas}

\vspace{0.3cm}

\ejercicio El conjunto $A = \{x \in \mathbb{N} \mid 2 < x \leq 6\}$
por extensi\'on es:
\begin{alternativas}
  \item $\{2, 3, 4, 5, 6\}$
  \item $\{3, 4, 5, 6\}$
  \item $\{2, 3, 4, 5\}$
  \item $\{3, 4, 5\}$
  \item $\{2, 4, 6\}$
\end{alternativas}

\vspace{0.3cm}

\ejercicio Si $A = \{a,\; b\}$, el n\'umero de subconjuntos de $A$ es:
\begin{alternativas}
  \item $2$
  \item $3$
  \item $4$
  \item $5$
  \item $6$
\end{alternativas}

\vspace{0.3cm}

\ejercicio Si $U = \{1,2,3,4,5\}$ y $A = \{2,4\}$, entonces $A' =$
\begin{alternativas}
  \item $\{1, 3, 5\}$
  \item $\{2, 4\}$
  \item $\{1, 2, 3\}$
  \item $\{3, 4, 5\}$
  \item $\emptyset$
\end{alternativas}

\end{multicols}

% ─── NIVEL INTERMEDIO ────────────────────────────────────────────────────────
\begin{tcolorbox}[
  enhanced, colback=AzulClaro!40, colframe=AzulMedio,
  arc=2pt, boxrule=0.5pt, fonttitle=\sffamily\bfseries\small,
  title={Nivel Intermedio}, left=4pt, right=4pt, top=3pt, bottom=3pt
]\end{tcolorbox}

\begin{multicols}{2}

\ejercicio Si $A = \{x \in \mathbb{Z}^{+} \mid \tfrac{24}{x} \in \mathbb{Z}\}$
y $B = \{y \in \mathbb{Z}^{+} \mid 18 = \overset{\circ}{y}\}$
(m\'ultiplos positivos de $y$ que dan 18), calcule $n(A) + n(B)$:
\begin{alternativas}
  \item $12$
  \item $14$
  \item $16$
  \item $18$
  \item $20$
\end{alternativas}

\vspace{0.3cm}

\ejercicio Si $n[\mathcal{P}(A)] = 32$, entonces el n\'umero de subconjuntos
propios de $A$ es:
\begin{alternativas}
  \item $16$
  \item $30$
  \item $31$
  \item $32$
  \item $64$
\end{alternativas}

\vspace{0.3cm}

\ejercicio Dados $A \subset B$, si $n[\mathcal{P}(B)] - n[\mathcal{P}(A)] = 112$,
entonces $n(B) + n(A)$ es:
\begin{alternativas}
  \item $9$
  \item $10$
  \item $11$
  \item $12$
  \item $13$
\end{alternativas}

\vspace{0.3cm}

\ejercicio Si $U = \{1,2,3,4,5,6\}$, $A = \{1,2,3,4\}$ y $B = \{3,4,5,6\}$,
entonces $n(A \mathbin{\triangle} B)$ es:
\begin{alternativas}
  \item $2$
  \item $3$
  \item $4$
  \item $5$
  \item $6$
\end{alternativas}

\vspace{0.3cm}

\ejercicio Dado $M = \{x \in \mathbb{Z} \mid x^{-x} = \tfrac{1}{4}\}$,
el n\'umero de subconjuntos propios de $M$ es:
\begin{alternativas}
  \item $1$
  \item $2$
  \item $3$
  \item $4$
  \item $7$
\end{alternativas}

\end{multicols}

% ─── NIVEL AVANZADO ──────────────────────────────────────────────────────────
\begin{tcolorbox}[
  enhanced, colback=NaranjaSuave!60, colframe=NaranjaTierra,
  arc=2pt, boxrule=0.5pt, fonttitle=\sffamily\bfseries\small,
  title={Nivel Avanzado --- Tipo Admisi\'on},
  left=4pt, right=4pt, top=3pt, bottom=3pt
]\end{tcolorbox}

\begin{multicols}{2}

\ejercicio Dado $U = \bigl\{\tfrac{x+3}{2} \in \mathbb{N}
\bigm| 5 < x < 17\bigr\}$, $D = \{d \in U \mid d \text{ es par}\}$
y $S = \{s \in U \mid s \leq 8\}$, calcule $n(D) + n(S)$:
\begin{alternativas}
  \item $7$
  \item $8$
  \item $9$
  \item $10$
  \item $11$
\end{alternativas}

\vspace{0.3cm}

\ejercicio Si $U = \{1,2,3\}$ y la proposici\'on
$\forall\, x,\; \forall\, y \;/\; x^{2} + y^{2} < 12$
es evaluada en $U \times U$, entonces es:
\begin{alternativas}
  \item Verdadera, pues $3^{2}+3^{2}=18>12$
  \item Falsa, pues $3^{2}+3^{2}=18>12$
  \item Verdadera, pues $1^{2}+1^{2}=2<12$
  \item Falsa, pues $1^{2}+1^{2}=2<12$
  \item Imposible determinarlo
\end{alternativas}

\vspace{0.3cm}

\ejercicio Dos conjuntos comparables tienen cardinales que difieren en $3$.
La suma de los cardinales de sus conjuntos potencia es $288$.
El cardinal del conjunto de menor tama\~no es:
\begin{alternativas}
  \item $3$
  \item $4$
  \item $5$
  \item $6$
  \item $7$
\end{alternativas}

\vspace{0.3cm}

\ejercicio Si $n(B-C)+n(C-B)=n(C)$, $n(B)=n(C)$,
$n[\mathcal{P}(A)]+n[\mathcal{P}(B)]=144$ y $n(A)=n(B)+3$,
calcula $n[\mathcal{P}(A)]+n[\mathcal{P}(B \cap C)]$:
\begin{alternativas}
  \item $136$
  \item $144$
  \item $152$
  \item $160$
  \item $168$
\end{alternativas}

\vspace{0.3cm}

\ejercicio Dados los conjuntos
$D = \{2x \mid x \in \mathbb{N},\; 0 < x < 6\}$,
$A = \bigl\{\tfrac{x+4}{2} \;\big|\; x \in D\bigr\}$ y
$N = \bigl\{\tfrac{2y+1}{3} \;\big|\; y \in A\bigr\}$,
calcule $n[\mathcal{P}(N)]$:
\begin{alternativas}
  \item $4$
  \item $8$
  \item $16$
  \item $32$
  \item $64$
\end{alternativas}

\end{multicols}

% ─── CLAVE DE RESPUESTAS ─────────────────────────────────────────────────────
\vspace{0.5cm}
\begin{tcolorbox}[
  enhanced, colback=GrisClaro, colframe=GrisMedio,
  arc=2pt, boxrule=0.5pt,
  fonttitle=\sffamily\bfseries\small\color{GrisCarbón},
  title={Clave de Respuestas}
]
\small\sffamily
\begin{multicols}{5}
\noindent
1.~E \quad 2.~D \quad 3.~B \quad 4.~C \quad 5.~A \\
6.~B \quad 7.~C \quad 8.~C \quad 9.~C \quad 10.~C \\
11.~C \quad 12.~B \quad 13.~C \quad 14.~A \quad 15.~B
\end{multicols}
\end{tcolorbox}

% FIN — ejercicios.tex
